\section{Introduction}\label{introduction}

\subsection{Context and motivation}\label{context-and-motivation}

The adoption of agents based on Large Language Models (LLMs) in software
engineering has been accompanied by new artifacts designed to
encapsulate operational knowledge. Among them are \emph{agent skills}:
portable directories that combine structured metadata, natural-language
instructions and, optionally, scripts and references. The
\texttt{SKILL.md} file is the descriptive core of the artifact and tells
the agent how to perform a task, which constraints to observe and which
resources to use \cite{agentskills_spec}. Skills extend the traditional
notion of reuse, since they distribute procedures and domain knowledge
without depending on a library with a formal interface.

Software reuse has been studied for decades
\cite{haefliger2008code,lopes2017dejavu}, but the diffusion of agent
skills has its own characteristics: there is no central registry or
package manager, and reuse happens by copying folders between
repositories. The same skill may appear in many repositories because it
is clear and portable, because it belongs to a template or catalog,
because a tool installed it, or because a single account duplicated it
across its own projects. These explanations mix properties of the
artifact and of the distribution context, which makes the observed reuse
hard to interpret.

GitSkills provides a large-scale corpus with skill occurrences, distinct
contents and repository metadata \cite{gitskills2027}. Other studies
examine specification defects \cite{zhang2026keepskillsreuse}, the
effect of skills on software engineering tasks
\cite{han2026sweskillsbench} and the evolution of prompts
\cite{tafreshipour2025prompting}. A complementary and methodologically
delicate question remains: which observable characteristics are
associated with the diffusion of a skill across distinct owners, and how
much of this predictability results from how the dataset itself was
built?

\subsection{Research problem}\label{research-problem}

We operationalize reuse as the presence of identical \texttt{SKILL.md}
content (same \texttt{content\_hash}) in distinct repositories. A first
version of this study modeled the number of repositories
(\texttt{repository\_count}) from skill attributes, text and attributes
of the representative repository, with repository-grouped validation and
an external test set. Two audits, described in this report, showed that
this design was fragile.

The first concerns the \textbf{representative}. GitSkills keeps detailed
metadata for only one copy of each content, chosen by the best location
class (preferring \texttt{.claude/skills/}) and then by the alphabetical
order of repository and path. Because a content with more copies has
more chances of having some copy in the preferred folder, the location
class of the representative, and the commit-history indicator that
depends on it, end up encoding the target itself. The second concerns
\textbf{self-copying and concentration}: 28\% of reused skills have all
copies under a single owner, and one account alone accounts for 12\% of
the development set. Counting repositories mixes diffusion among
developers with duplication by a single account.

We therefore reformulated the problem. The main target becomes
\textbf{cross-owner reuse} (\texttt{cross\_owner}: the content appears
under two or more distinct owners) and, secondarily, the number of
owners (\texttt{owner\_count}). The attributes now exclude the variables
that depend on the representative rule, which are analyzed as a finding.
The evaluation uses owner-grouped validation, a per-owner cap, nested
tuning and pre-specified hypotheses, and confirmation is done on a
\textbf{new holdout, disjoint by owner}, evaluated once. Because the
study was revised after exploratory results, we explicitly distinguish
what is confirmatory from what is exploratory. Conclusions are
observational associations \cite{shmueli2010explain}: we claim neither
causality nor skill quality.

\subsection{Objectives and research
questions}\label{objectives-and-research-questions}

The general objective is to investigate, reproducibly, which observable
characteristics are associated with the reuse of agent skills across
owners and how much this association is affected by dataset construction
artifacts.

\begin{itemize}
\tightlist
\item
  \textbf{RQ1 --- Construction biases:} to what extent do the
  representative-selection rule, self-copying and owner concentration
  confound the target and the attributes of GitSkills?
\item
  \textbf{RQ2 --- Predictability:} with attributes free of these
  artifacts, can cross-owner reuse and the number of owners be predicted
  better than baselines, and is the estimate confirmed on an
  owner-disjoint holdout?
\item
  \textbf{RQ3 --- Origin of the signal:} what is the contribution of
  skill, repository and catalog attributes, which variables are
  associated with the predictions, and how consistently?
\end{itemize}

\subsection{Contributions}\label{contributions}

\begin{enumerate}
\def\labelenumi{\arabic{enumi}.}
\tightlist
\item
  An \textbf{audit of the GitSkills representative}: we recover its rule
  (97.2\% agreement) and show that location and history depend on the
  number of copies.
\item
  A \textbf{protocol} for modeling cross-owner reuse, with an
  owner-based target, a per-owner cap, grouped validation, pre-specified
  hypotheses with Holm correction and confirmation on a new holdout.
\item
  Evidence of a \textbf{modest predictive signal}, with a stable RMSLE
  gain across protocol variations, decomposed into blocks (skill,
  repository, catalog, path) and confirmed on unseen data: AUC 0.73 for
  cross-owner reuse.
\item
  A \textbf{reproducible package} that records the data revision, seeds,
  partitions and a lock that prevents reusing the holdout.
\end{enumerate}

\subsection{Organization}\label{organization}

Section 2 presents related work; Section 3, the methodology; Section 4,
the results; Section 5, the discussion and threats to validity; Section
6, the conclusions. The appendices cover reproduction and the use of AI
tools.
